\bookchapter{第六章：守護汝之數字生活——安全與隱私}

\chapterintro{手機之安=隱私之安加財產之安。於數據即資產之時，護手機即護汝之己。}

\booksection{一、鎖屏：第一道物理防線}

鎖屏乃汝手機之第一道防線，亦乃最易被忽者。

\booksubsection{安之級之序：}

\begin{enumerate}
  \item \textbf{密碼/PIN碼}：安度至高，乃最終之保（忘指紋/臉時賴之）。
  \item \textbf{指紋之識}：安與便之完之衡。
  \item \textbf{人臉之識}：至便，然安之性賴於技（3D結構光安，2D照片易騙）。
  \item \textbf{圖案之解}：便而易留痕被人猜。
\end{enumerate}

\booksubsection{設強鎖屏：}

\begin{itemize}
  \item 務必設一複雜之密碼（勿用生日、123456）。
  \item 開生物之識。
  \item 開生物之識。
\end{itemize}

\booksection{二、“查找吾手機”：終極防盜之具}

\textbf{必開！}其能：

\begin{itemize}
  \item \textbf{定位}：手機失於何，於圖上尋之。
  \item \textbf{放聲}：使手機響鈴，宜家中尋時。
  \item \textbf{丟失之模}：遠鎖死手機並顯聯之號。
  \item \textbf{遠擦除}：最後一招，清所有數據（然擦後不能再定位）。
\end{itemize}

往 [設置]，蘋果尋“查找吾iPhone”，安卓尋“查找吾設備”或品牌類之服，保其爲綠開之態。

\booksection{三、系統之更：打補丁，堵洞}

系統之更首旨在修安之洞，猶爲屋補牆洞。故，見系統更之提，盡即時升之。

系統之更首旨在修安之洞，猶爲屋補牆洞。故，見系統更之提，盡即時升之。

\booksection{四、識與防電信詐}

於AI技迅猛之今，傳統之“眼見爲實”之識之模已全失效。詐者正用深僞（Deepfake）等AI技，實現自“廣撒網”式詐至“精圍獵”式罪之全升。

\booksubsection{AI 時代最常之數詐之套}

\textbf{(1) AI 換臉擬聲之詐}\par

(1) AI 換臉擬聲之詐\textbf{突車禍、緊借、遭綁}等。此類之詐用親情、友情等情之紐，\textbf{迷之性極強}。

\textbf{。}\par

(2) 冒權威/客服之詐\textbf{造恐慌之心或權威之壓}，誘受害者轉或開屏共之功，從而竊賬戶之信。

\textbf{(3) AI 投資理財之詐}\par

(3) AI 投資理財之詐\textbf{終捲款而逃}。

\booksubsection{破局之法——“照妖鏡”防護之思}

AI能克聲與面，然亦非無懈，有數破可利用：

\begin{itemize}
  \item \textbf{動作驗證法}：行視頻通話時，主動要對方完預設之連貫動作之令，如"用手於臉前速揮過"、"側臉示耳之廓"或"連速眨五次"。今之AI換臉技於處複雜動態之遮、速動之景與光影之變時，\textbf{極易現畫之卡頓、面之扭或邊之穿模}等破。
  \item \textbf{暗號覈驗法}：建並用預設之私核問，如"吾上次同聚時汝點何菜"、"吾送汝之首物何"等惟雙方知之細。若對方顯遲疑或答非所問，應\textbf{即止通話並啓安之預}。
  \item \textbf{二次覈實法}：嚴守公安部網安局所發之"\textbf{四不一核}"之則——不輕信陌之信、不隨轉匯、不泄個人敏信、不傳可疑之內容。凡涉錢財往來，務必經\textbf{已知之原號回撥或當面認}。行轉之事時，建議擇"二十四時到"之模，爲潛詐之識與資之挽預足時窗。
\end{itemize}

\booksubsection{傳統常之套（仍須警）：}

\begin{itemize}
  \item \textbf{中獎/刷單返利}：利誘，先予小甜再騙大錢。
  \item \textbf{網戀之詐}：殺豬盤，養情後騙錢。
  \item \textbf{冒熟人}：盜號後借（AI 時代更難辨）。
\end{itemize}

\booksubsection{手機自帶之護：}

\begin{itemize}
  \item 騷擾電話之攔。
  \item 驗證碼防竊。
  \item 驗證碼防竊。
\end{itemize}

\booksubsection{防騙之訣：}

\textbf{不輕信、不泄、不點、不轉、有疑問先核！}遇“急事”先冷五分，幾所有詐於此時將露馬腳。

\booksection{五、御“開盒”與社工攻：斬隱私泄之鏈}

“開盒”（社會工程學攻）之核之理在\textbf{系統性集、析並拼個人數字之片，構建完之個人之像}。防之要在爲個人隱私劃明紅線，建多層次之護之體。

\booksubsection{照之脫敏處}

基礎 EXIF 清理方法見第二章“圖片與文件系統管理”。於社交媒體分享照片前，尚需進一步\textbf{徹刪照之元數據（EXIF信）}，含攝之時、設備之型、GPS位之座標等敏之數據。於照中所含之\textbf{門牌號、車牌之信、車票之條碼、身份證號}等關鍵之信嚴打碼或模糊處，避無中泄個人行蹤之軌。

\begin{itemize}
  \item iOS：[設置] > [隱私與安全性] > [相機] 中可閉"精確位置"；分享時擇"僅當前照"可剝元數據。
  \item 安卓：部分相冊App（如Google相冊）分享時將提示是否保留位置之信。
  \item 通用之具：ExifTool、Metapho 等可查/編 EXIF。
\end{itemize}

\booksubsection{應用權限之收}

基礎權限管理方法詳見第四章“權限管理”。於“防開盒”場景下，尚需進一步定期審查並關閉電商類、工具類App之\textbf{麥克風、通訊錄、位之信}等非必之權，防App於後臺“偷聞”日之對或非法集個人關係之網；對久置而猶佔高權之應用宜即卸，從源頭減隱私泄之險點。

\begin{itemize}
  \item 見"通訊錄之權"請時，先問之：此App果必知吾所有聯人乎？
  \item 三月未用之App，果卸。權限乃越用越緊，非越用越松。
\end{itemize}

\booksubsection{智能設備安之設防}

智能攝頭、智能門鎖、智能音箱等IoT設備於初用時必\textbf{即改默認登之密}，設含大小寫字母、數字與特符之高強獨立之密；不常用之遠訪功宜常閉，定期檢設備固件之更並即補安之洞，\textbf{防黑客侵併窺家隱私之空}。

\begin{tcbwarn}{真案之警}
真案之警\textbf{不改默認密 = 將家大門敞}。

\end{tcbwarn}

\booksection{六、隱私之護：汝之數據乃汝之財}

手機中之照、聊天記、位之軌，本質上皆汝之“數字之資”。一泄，輕則告擾，重則人肉、詐甚至敲詐。

\booksubsection{App 權限之“最小化之則”：}

\begin{itemize}
  \item \textbf{位置}：打車/地圖用"始終允"；他App用"僅用時允"。
  \item \textbf{通訊錄}：除非必（如微信加好友）方予，否則一律禁。
  \item \textbf{相冊/全部照}：iOS 14+ 起可授"選定之照"，安卓 13+ 同支持精授。
  \item \textbf{麥克風/攝頭}：平閉，需時再開。疑被偷錄時，滑出控制中心觀有無"小綠點/小橙點"。
  \item \textbf{後臺之跟}：iOS 14.5+ 起強要App用 IDFA 前必問；安卓可閉"廣告ID"。
\end{itemize}

\booksubsection{隱私之管實戰：}

\begin{itemize}
  \item \textbf{隱私之報（iOS）}：[設置] > [隱私與安全性] > [App 隱私之報]，觀何App於偷訪汝之位、相機、麥克風。
  \item \textbf{權限用之記（安卓）}：[設置] > [隱私] > [權限用之記]，觀誰於何時調了敏之權。
  \item \textbf{限廣告之跟}：於系統設置中閉"廣告ID"或"限廣告之跟"，能減"被精廣告擾"。
\end{itemize}

\booksection{七、失/被盜後之“急救”單}

手機失勿慌，依序處之：

\begin{enumerate}
  \item \textbf{即凍 SIM 卡}：聯運營商掛失，避被用以收驗證碼。
  \item \textbf{遠鎖}：經雲服（iCloud/品牌雲）開"丟失之模"。
  \item \textbf{改要之密}：登過手機之所有要之賬號（支付、銀行、郵箱）即改密。
  \item \textbf{報案}：往派出所報案，錄 IMEI 號。手機尋回時亦便於證身份。
  \item \textbf{告親友}：防騙以汝名借錢。
  \item \textbf{（最後）遠擦}：確認尋無望時再操，擦後定位之功失效。
\end{enumerate}

\booksection{八、SIM 卡與 eSIM 之安}

\begin{itemize}
  \item \textbf{設 SIM 卡 PIN}：開後，每重啓手機或換卡皆須輸 PIN碼，防"拔卡換手機"以繞鎖屏。
  \item \textbf{eSIM 之優}：無實體卡可被拔，安之性更高；同時支持多號切，旅時甚爲便。
  \item \textbf{二次換卡之險}：警"運營商二次換卡"之騙——騙者冒用汝之身補辦 SIM 卡，從而接管汝之手機號。設"補卡須機主本人持身份證至營業廳辦"乃最效之防。
\end{itemize}

\begin{tcbwarn}{切勿爲之事}
\begin{itemize}
  \item 用生日/123456/qwerty 等弱密；
  \item 將驗證碼予任何人（縱爲"客服"）；
  \item 將驗證碼予任何人（縱爲"客服"）；
  \item 於公共 Wi-Fi 下登網銀或支付；
  \item 予陌人開"屏共"或"遠控"。
\end{itemize}

\end{tcbwarn}

\begin{tcbpractice}{試啓並擬"遠鎖"}
\begin{enumerate}
  \item 保汝"查找吾手機"之功已開。
  \item 保汝"查找吾手機"之功已開。
  \item 尋"查找設備"之功，試對汝己之手機發一"放聲"或擬"鎖"之令，熟全操之流。
\end{enumerate}

\end{tcbpractice}
